\documentclass{amsart}

% 数学相关宏包
\usepackage{amsmath,mathtools,amsthm,amssymb}
\usepackage{bbm}
\usepackage{mathrsfs}
\usepackage{centernot}

% 图形和绘图
\usepackage{graphicx,float}
\usepackage{tikz,tikz-cd}
\usepackage{pgfplots}
\usepackage[framemethod=TikZ]{mdframed}
\usetikzlibrary{decorations.pathmorphing,positioning,arrows,automata,calc,shapes,patterns}
\pgfplotsset{compat=1.18}
\usepgfplotslibrary{statistics}

% 表格和列表
\usepackage{booktabs,multirow,colortbl}
\usepackage{enumitem}
\setlist[itemize,enumerate,description]{noitemsep}
\setlistdepth{9}

% 算法
\usepackage{algorithm,algorithmic}

% 字体和颜色
\usepackage{xcolor,listings}
\usepackage{xeCJK}
\setCJKmainfont{Noto Serif CJK SC}

% 页面和链接
\usepackage{fancyhdr,caption}
\usepackage{hyperref}
\usepackage{cleveref}

% 工具宏包
\usepackage{etoolbox,letltxmacro,docmute}
\usepackage{gensymb}

% 自定义设置
\renewcommand{\arraystretch}{1.5}
\let\originalmathbb\mathbb
\renewcommand{\mathbb}[1]{%
    \ifstrequal{#1}{1}{\mathbbm{#1}}{\originalmathbb{#1}}%
}

% 颜色定义
\definecolor{lightblue}{RGB}{173,216,230}
\definecolor{lightgreen}{RGB}{144,238,144}
\definecolor{lightyellow}{RGB}{255,255,224}
\definecolor{lightgray}{RGB}{211,211,211}

		% 超链接设置
		\hypersetup{
		    colorlinks=true,
		    linkcolor=red,
		    filecolor=magenta,
		    urlcolor=cyan,
		    pdftitle={\@title},
		    pdfauthor={\@author},
		    pdfsubject={数学笔记},
		    pdfkeywords={数学},
		    bookmarksnumbered=true,
		    bookmarksopen=true,
		    bookmarksopenlevel=6,
		    pdfstartview=FitH,
		    pdfpagemode=UseOutlines,
		    pdfborder={0 0 0}
		}

% 定理环境
\theoremstyle{plain}
\newtheorem{theorem}{Theorem}[section]
\newtheorem{lemma}[theorem]{Lemma}
\newtheorem{corollary}[theorem]{Corollary}
\newtheorem{proposition}[theorem]{Proposition}

\theoremstyle{definition}
\newtheorem{definition}[theorem]{Definition}
\newtheorem{example}[theorem]{Example}
\newtheorem{exercise}[theorem]{Exercise}
\newtheorem{problem}[theorem]{Problem}

\theoremstyle{remark}
\newtheorem{remark}[theorem]{Remark}
\newtheorem{note}[theorem]{Note}
\newtheorem{observation}[theorem]{Observation}

% 解答环境
\newtheorem*{solution}{Solution}
\newtheorem*{proof*}{Proof}

% 图片路径
\graphicspath{{F:/iCloudDrive/iCloud~md~obsidian/2025-summer/figures/}}

\author{乐绎华, Yue Yihua}
\address{中山大学, Sun Yat-sen University}
\email{yueyh@mail2.sysu.edu.cn}
\date{\today}
\title{Mathematical Notes}

\begin{document}

\maketitle
\tableofcontents

See 随机过程基础 Basic Stochastic Processes (Zdzisław Brzeźniak, Tomasz Zastawniak) (Z-Library).

Bachelier assumed the relative infinitesimal price increments $dX(t)$ are proportional to the increments $dW(t)$ of the Wiener process,
\[
\frac{dX(t)}{X(t)} =\sigma dW(t)\iff X(t)=X(0)+\sigma \int_{0}^{t} X(t) \, \mathrm{d}W(t)  
\]
where $\sigma$ is a positive constant, the right hand side is called the Ito integral. The solution is expected to be $xe^{ W(t) }$, but in fact
\[
X(t)=xe^{ W(t) }e^{ -\frac{t}{2} }
\]
due to the non-differentiability of the paths of the Wiener process.

We define Ito integral from the same spirit of Riemann integral, but the convergence is in $L^{2}$ instead of $\mathbb{R}$, and Riemann integral
\[
\sum_{i=1}^{n} f(s_i)(x_i-x_{i-1})
\]
does not depend on $s_i\in[x_{i-1},x_i]$, but in stochastic case,
\[
\sum_{i=1}^{n} f(s_i)(W(x_i)-W(x_{i-1}))
\]
it does.

\begin{figure}[H]
\centering
\includegraphics[width=\textwidth]{Pasted image 20250723145517.png}
% \caption{}
\label{}
\end{figure}
\begin{proof}
\[
\begin{aligned}
\sum W(t_j^{n})(W(t_{j+1}^{n})-W(t_j^{n})) & =\sum \frac{1}{2}(W(t_{j+1}^{n})^{2}-W(t_j^{n})^{2})-\frac{1}{2}\underbrace{ (W(t_j^{n})-W(t_{j+1}^{n}))^{2} }_{ =(\Delta _i^{n}W)^{2} } \\
 & =\frac{1}{2}W(t_{n}^{n})^{2}-\frac{1}{2}W(t_0^{n})^{2}-\frac{1}{2}\sum(\Delta _i^{n}W)^{2} \\
 & =\frac{1}{2}(W(T)^{2}-W(0)^{2})-\frac{1}{2}\sum(\Delta _i^{n}W)^{2} \\
 & \to\frac{1}{2}(W(T)^{2}-W(0)^{2})-\frac{1}{2}T
\end{aligned}
\]
\[
\begin{aligned}
\sum W(t_{j+1}^{n})(W(t_{j+1}^{n})-W(t_j^{n})) & =\sum \frac{1}{2}(W(t_{j+1}^{n})^{2}-W(t_j^{n})^{2})+\frac{1}{2}\underbrace{ (W(t_j^{n})-W(t_{j+1}^{n}))^{2} }_{ =(\Delta _i^{n}W)^{2} }  \\
 & \to\frac{1}{2}(W(T)^{2}-W(0)^{2})+\frac{1}{2}T
\end{aligned}
\]
\end{proof}

\begin{figure}[H]
\centering
\includegraphics[width=\textwidth]{1-ito-stochastic-calculus-2025072314.png}
% \caption{}
\label{}
\end{figure}
\begin{figure}[H]
\centering
\includegraphics[width=\textwidth]{2-ito-stochastic-calculus-2025072314.png}
% \caption{}
\label{}
\end{figure}
\begin{figure}[H]
\centering
\includegraphics[width=\textwidth]{ito-stochastic-calculus-2025072314.png}
% \caption{}
\label{}
\end{figure}

The choice of $s_j$ makes difference. To ensure that the process adapted to $\mathcal{F}_{t}$, we have to take $s_j=t_j$.

\begin{figure}[H]
\centering
\includegraphics[width=\textwidth]{ito-stochastic-calculus-2025072315.png}
% \caption{}
\label{}
\end{figure}
\begin{figure}[H]
\centering
\includegraphics[width=\textwidth]{1-ito-stochastic-calculus-2025072315.png}
% \caption{}
\label{}
\end{figure}

\begin{figure}[H]
\centering
\includegraphics[width=\textwidth]{2-ito-stochastic-calculus-2025072315.png}
% \caption{}
\label{}
\end{figure}
\begin{figure}[H]
\centering
\includegraphics[width=\textwidth]{3-ito-stochastic-calculus-2025072315.png}
% \caption{}
\label{}
\end{figure}
\begin{proof}
Expand both sides.
\end{proof}

\begin{figure}[H]
\centering
\includegraphics[width=\textwidth]{4-ito-stochastic-calculus-2025072315.png}
% \caption{}
\label{}
\end{figure}
\begin{proof}
Trivial.
\end{proof}

\begin{figure}[H]
\centering
\includegraphics[width=\textwidth]{5-ito-stochastic-calculus-2025072315.png}
% \caption{}
\label{}
\end{figure}
\begin{proof}
Trivial.
\end{proof}
\begin{figure}[H]
\centering
\includegraphics[width=\textwidth]{6-ito-stochastic-calculus-2025072315.png}
% \caption{}
\label{}
\end{figure}

\begin{figure}[H]
\centering
\includegraphics[width=\textwidth]{7-ito-stochastic-calculus-2025072315.png}
% \caption{}
\label{}
\end{figure}

The Ito stochastic integral is defined by approximation.

\begin{figure}[H]
\centering
\includegraphics[width=\textwidth]{8-ito-stochastic-calculus-2025072315.png}
% \caption{}
\label{}
\end{figure}
\begin{proof}
Prove by approximation in $L^{2}$ and $M^{2}$.
\end{proof}

\begin{figure}[H]
\centering
\includegraphics[width=\textwidth]{9-ito-stochastic-calculus-2025072315.png}
% \caption{}
\label{}
\end{figure}

We define the Ito integral over $[0,T]$.

\begin{figure}[H]
\centering
\includegraphics[width=\textwidth]{10-ito-stochastic-calculus-2025072315.png}
% \caption{}
\label{}
\end{figure}

\begin{figure}[H]
\centering
\includegraphics[width=\textwidth]{12-ito-stochastic-calculus-2025072315.png}
% \caption{}
\label{}
\end{figure}

\begin{proof}
\begin{figure}[H]
\centering
\includegraphics[width=\textwidth]{13-ito-stochastic-calculus-2025072315.png}
% \caption{}
\label{}
\end{figure}
\begin{figure}[H]
\centering
\includegraphics[width=\textwidth]{14-ito-stochastic-calculus-2025072315.png}
% \caption{}
\label{}
\end{figure}
\end{proof}

To verify the existence of integral, it's not always easy to find the approximation of step functions. Similar to Riemann integral, "continuous functions are always integrable", there is a corresponding argument for Ito integral.
\begin{figure}[H]
\centering
\includegraphics[width=\textwidth]{15-ito-stochastic-calculus-2025072315.png}
% \caption{}
\label{}
\end{figure}

\begin{proof}
Omitted.
\end{proof}

\begin{figure}[H]
\centering
\includegraphics[width=\textwidth]{16-ito-stochastic-calculus-2025072315.png}
% \caption{}
\label{}
\end{figure}

\begin{figure}[H]
\centering
\includegraphics[width=\textwidth]{17-ito-stochastic-calculus-2025072315.png}
% \caption{}
\label{}
\end{figure}
\begin{figure}[H]
\centering
\includegraphics[width=\textwidth]{18-ito-stochastic-calculus-2025072315.png}
% \caption{}
\label{}
\end{figure}

\begin{figure}[H]
\centering
\includegraphics[width=\textwidth]{20-ito-stochastic-calculus-2025072315.png}
% \caption{}
\label{}
\end{figure}
\begin{figure}[H]
\centering
\includegraphics[width=\textwidth]{1-ito-stochastic-calculus-2025072316.png}
% \caption{}
\label{}
\end{figure}

\begin{figure}[H]
\centering
\includegraphics[width=\textwidth]{ito-stochastic-calculus-2025072316.png}
% \caption{}
\label{}
\end{figure}
\begin{proof}
Use Jensen's formula.
\end{proof}
\begin{figure}[H]
\centering
\includegraphics[width=\textwidth]{2-ito-stochastic-calculus-2025072316.png}
% \caption{}
\label{}
\end{figure}

Next we consider the stochastic differentiation, which is different from the classical calculus for smooth function.

\begin{figure}[H]
\centering
\includegraphics[width=\textwidth]{3-ito-stochastic-calculus-2025072316.png}
% \caption{}
\label{}
\end{figure}
\begin{figure}[H]
\centering
\includegraphics[width=\textwidth]{4-ito-stochastic-calculus-2025072316.png}
% \caption{}
\label{}
\end{figure}

\begin{figure}[H]
\centering
\includegraphics[width=\textwidth]{5-ito-stochastic-calculus-2025072316.png}
% \caption{}
\label{}
\end{figure}
\begin{figure}[H]
\centering
\includegraphics[width=\textwidth]{6-ito-stochastic-calculus-2025072316.png}
% \caption{}
\label{}
\end{figure}
\begin{figure}[H]
\centering
\includegraphics[width=\textwidth]{7-ito-stochastic-calculus-2025072316.png}
% \caption{}
\label{}
\end{figure}
\[
F_{t}'(t,x)=-\frac{1}{2}e^{ x-\frac{t}{2} },\quad F_{xx}''(t,x)=e^{ x-\frac{t}{2} },\quad F_{x}'(t,x)=e^{ x-\frac{t}{2} }
\]
are continnuous for all $t\geq0$ and $x\in \mathbb{R}$. Assume that $F_{x}' (t,W (t))\in M_{T}^{2}$ for all $T\geq0$. Then $F(t,W(t))$ is Ito process and
\[
dF(t,W(t)) = \underbrace{ \left( -\frac{1}{2}e^{ W(t) }e^{ -\frac{t}{2} }+\frac{1}{2}e^{ W(t) }e^{ -\frac{t}{2} } \right) }_{ =0 }dt+\underbrace{ e^{ W(t) }e^{ -\frac{t}{2} } }_{ =F(t,W(t)) }dW(t) 
\]
i.e.
\[
dX(t)=X(t)dW(t)
\]
\begin{figure}[H]
\centering
\includegraphics[width=\textwidth]{8-ito-stochastic-calculus-2025072316.png}
% \caption{}
\label{}
\end{figure}

Then we consider the Stochastic Differential Equations.
\[
d\xi(t)=f(\xi(t))dt+g(\xi(t))dW(t)
\]
\begin{figure}[H]
\centering
\includegraphics[width=\textwidth]{9-ito-stochastic-calculus-2025072316.png}
% \caption{}
\label{}
\end{figure}
\begin{figure}[H]
\centering
\includegraphics[width=\textwidth]{10-ito-stochastic-calculus-2025072316.png}
% \caption{}
\label{}
\end{figure}

\begin{figure}[H]
\centering
\includegraphics[width=\textwidth]{11-ito-stochastic-calculus-2025072316.png}
% \caption{}
\label{}
\end{figure}
\begin{figure}[H]
\centering
\includegraphics[width=\textwidth]{12-ito-stochastic-calculus-2025072316.png}
% \caption{}
\label{}
\end{figure}

We have Existence and Uniqueness theorem.

\begin{figure}[H]
\centering
\includegraphics[width=\textwidth]{13-ito-stochastic-calculus-2025072316.png}
% \caption{}
\label{}
\end{figure}


\end{document}